% =================================================================
% Deep Hedging Under Market Frictions, Crisis Regimes, and Exotic Payoffs
% MSc / Thesis -- Pratham Kailasiya, IIT Roorkee
% =================================================================
\documentclass[11pt,a4paper]{article}

% ---------- Packages ---------------------------------------------------------
\usepackage[a4paper,margin=2.4cm]{geometry}
\usepackage[T1]{fontenc}
\usepackage{lmodern}
\usepackage{amsmath,amssymb,amsthm}
\usepackage{mathtools}
\usepackage{booktabs}
\usepackage{tabularx}
\usepackage{microtype}
\usepackage{enumitem}
\usepackage{graphicx}
\usepackage{caption,subcaption}
\usepackage{hyperref}
\usepackage{algorithm,algorithmic}
\usepackage{xcolor}
\usepackage{titlesec}
\usepackage[round,authoryear]{natbib}
\bibliographystyle{plainnat}

\hypersetup{colorlinks=true,linkcolor=blue!50!black,citecolor=blue!50!black,
            urlcolor=blue!50!black}

% ---------- Math shortcuts ---------------------------------------------------
\newcommand{\R}{\mathbb{R}}
\newcommand{\E}{\mathbb{E}}
\newcommand{\PP}{\mathbb{P}}
\newcommand{\QQ}{\mathbb{Q}}
\newcommand{\F}{\mathcal{F}}
\newcommand{\cvar}{\mathrm{CVaR}}
\newcommand{\VaR}{\mathrm{VaR}}
\DeclareMathOperator*{\argmin}{arg\,min}
\DeclareMathOperator*{\argmax}{arg\,max}

\newtheorem{theorem}{Theorem}
\newtheorem{proposition}[theorem]{Proposition}
\newtheorem{lemma}[theorem]{Lemma}
\newtheorem{definition}{Definition}
\newtheorem{corollary}[theorem]{Corollary}
\newtheorem*{remark}{Remark}
\newtheorem{assumption}{Assumption}

\titleformat{\section}[hang]{\large\bfseries}{\thesection}{0.6em}{}

% =============================================================================
\title{\textbf{Deep Hedging under Market Frictions, Crisis Regimes,\\ and
Exotic Payoffs:\\\large An Empirical Study Across Vanilla, Barrier, Digital,
Asian and Basket Options on Synthetic and Real Markets}}

\author{Pratham Kailasiya\\
Department of Mathematics, IIT Roorkee\\
\texttt{\small prathamkailasiya@gmail.com}}

\date{\today}

\begin{document}
\maketitle

\begin{abstract}
\noindent
Deep hedging \citep{buehler2019deep} learns option-replicating policies
under market frictions, but most published evaluations are restricted to a
single calibrated Heston regime, vanilla European options, and a single
underlier.  This thesis closes that gap on three axes simultaneously:
(i)~\emph{regime}: Heston, multi-asset Bates jump-diffusion, six historical
crisis windows on real US and Indian equity feeds; (ii)~\emph{payoff}:
European call/put, knock-out / knock-in barrier, cash-or-nothing digital,
arithmetic Asian, and three-asset equally-weighted basket calls;
(iii)~\emph{architecture}: LSTM, causal Transformer, Wasserstein-DRO
Transformer (W-DRO-T), a three-stage curriculum LSTM (3SCH), and a
regime-switching ensemble (RSE) over LSTM/Transformer/Signature experts.
We formalise the optimisation problems, prove the relevant duality results,
audit and reimplement all five hedgers under one fair training protocol, and
compare them on $7$ tasks $\times$ $6$ regimes $\times$ $5$ random seeds with
bootstrap confidence intervals and Holm--Bonferroni-corrected paired tests.
\textbf{Headline result}: under a unified protocol the causal Transformer
and W-DRO-T dominate the LSTM benchmark by $13$--$30\%$ in CVaR$_{95}$
across vanilla, barrier and Asian payoffs in normal and COVID regimes; W-DRO-T
yields the largest tail-risk reduction in stress regimes and the lowest
P\&L volatility, while RSE attains the lowest coefficient of variation
across seeds.  3SCH dominates only in the multi-asset basket-call setting
where its low turnover offsets its higher tail risk under high-friction
markets.  We supply a fully reproducible pipeline, bootstrap-validated
benchmark tables, ablation studies, sensitivity analyses, and a thorough
discussion of economic, regulatory and managerial implications.
\end{abstract}

\tableofcontents
\newpage

% =============================================================================
\section{Introduction}\label{sec:intro}

\subsection{Motivation}

Hedging an option position is the fundamental risk-management problem in
quantitative finance.  Black--Scholes--Merton theory \citep{black1973pricing,
merton1973theory} prescribes continuous delta hedging at zero cost; in practice,
trading is discrete, transaction costs are non-zero, volatility is
stochastic, and the distribution of returns shifts dramatically between
regimes.  These deviations from the idealised model break dynamic
replication and produce non-trivial residual P\&L.  \emph{Deep hedging}
\citep{buehler2019deep} reformulates the problem as a stochastic optimal
control problem in which a neural network is trained directly to minimise a
convex risk measure of the residual hedging error; the resulting policies
account for frictions and non-linear payoffs without analytical pricing
formulae.

The empirical literature has so far focused on three narrow but tractable
choices: a single calibrated Heston model, a single European call, and a
single equity underlier.  Real desks face a wider problem: they hedge
exotics whose payoffs depend on path geometry; they hold portfolios with
several correlated underliers; and the calm/stress regime in which they
operate cycles unpredictably between calm and crisis (the GFC, COVID and
banking-stress episodes are now part of the empirical record).  An
honest assessment of which architecture should be deployed must therefore
test models \emph{across} regime, payoff and asset count simultaneously.

\subsection{Research questions}

\begin{enumerate}[label=\textbf{Q\arabic*.},leftmargin=*]
\item Do learned policies that win on a single Heston regime continue to win
  in multi-asset Bates jump-diffusion environments and on real
  market data sliced into historical crises?
\item Does the relative ranking of architectures (LSTM vs. Transformer vs.
  W-DRO-T vs.\ 3SCH vs.\ RSE) depend on payoff (vanilla, barrier, digital,
  Asian, basket) and on the friction level?
\item What is the marginal value of each design choice (DRO penalty,
  curriculum stages, ensemble gating) when measured under a uniform
  training budget and against a strong recurrent baseline?
\item What economic, regulatory and managerial implications follow?
\end{enumerate}

\subsection{Contributions}

\begin{enumerate}[leftmargin=*]
\item \textbf{An extended evaluation suite.}  We audit, refactor, and run
  five hedgers on a $7$-task suite (European call/put, up-and-out barrier,
  down-and-in barrier, digital, Asian, basket) under six regimes
  (normal/post-COVID/COVID/GFC US, normal/COVID India) and a multi-asset
  Bates jump-diffusion environment.

\item \textbf{Online + walk-forward real-data validation.}  A reproducible
  pipeline that pulls SPY, QQQ, NIFTY data via \texttt{yfinance}, slices it
  into crisis windows (GFC 2008, vol-Q4-2018, COVID 2020, inflation 2022,
  banking 2023), and runs walk-forward train/test on hedging episodes
  cut from the price series.  An offline-safe synthetic fallback ensures
  reproducibility without internet access.

\item \textbf{Tail-risk benchmark with statistical guarantees.}  All
  comparisons use $5$ independent seeds with bootstrap-resampled $95\%$
  CIs ($2{,}000$ resamples), paired $t$-tests against the LSTM baseline,
  and Holm--Bonferroni multiple-test correction.

\item \textbf{Honest ranking of existing methods.}  We find Transformer-based
  policies (vanilla and W-DRO-T) consistently beat the LSTM baseline; the
  three-stage curriculum 3SCH offers no statistically significant tail-risk
  improvement on single-asset tasks, but does dominate on the multi-asset
  basket where its low turnover is decisive.  RSE has the smallest
  cross-seed variance, which matters in production.

\item \textbf{Reproducible artefacts.}  Code, fitted models, pickled
  metrics, raw P\&L distributions, and figure-generating scripts are
  released alongside this thesis.
\end{enumerate}

% =============================================================================
\section{Background and Literature Review}\label{sec:lit}

\subsection{Classical replication}

Black--Scholes \citep{black1973pricing} and Merton \citep{merton1973theory}
established that, under continuous trading, log-normal dynamics, and zero
transaction costs, the European call admits a unique replicating strategy
$\delta_t^{\mathrm{BS}} = \Phi(d_1)$.  Real markets violate every
assumption: \citet{leland1985option} showed that proportional transaction
costs require an inflated effective volatility; \citet{whalley1997asymptotic}
derived the optimal no-trade band under exponential utility for small costs.
These analytical adjustments are tied to specific dynamics and payoff types
and break down under jumps, stochastic volatility, or path-dependence.

Stochastic-volatility models extend the underlying dynamics:
\citet{heston1993closed} proposed an affine variance process admitting
semi-closed pricing formulas; \citet{bates1996jumps} added i.i.d.\ Poisson
jumps; \citet{barndorff2002econometric} formalised realised-volatility
estimators.  Markets are incomplete under stochastic volatility and so
\emph{any} hedging strategy leaves residual risk; the question becomes
\emph{which} risk measure to optimise.  Coherent risk measures
\citep{artzner1999coherent} and convex risk measures
\citep{follmer2002convex} provide the axiomatic backbone; CVaR has the
useful Rockafellar--Uryasev dual \citep{rockafellar2000optimization} that
makes it amenable to gradient-based optimisation.

\subsection{Deep hedging}

\citet{buehler2019deep} parameterise the hedging policy as a neural
network and minimise an empirical convex risk measure across simulated
paths.  They show that the learned policy outperforms Black--Scholes in
the presence of transaction costs.  Several extensions are now standard:
\citet{kozyra2019deep} introduces a two-stage curriculum (CVaR pretraining
followed by entropic fine-tuning) and the $\tanh$ delta clamp;
\citet{horvath2021deep} extends to rough volatility;
\citet{carbonneau2021deep} applies the framework to long-dated insurance
liabilities; \citet{wiese2020deep} learns the underlying dynamics jointly.
\citet{lutkebohmert2021robust} and the W-DRO line \citep{blanchet2019quantifying,
esfahani2018data,rahimian2019distributionally} introduce distributionally
robust hedging via Wasserstein ambiguity sets.  \citet{tong2023sigformer}
combines path signatures \citep{lyons1998differential,kidger2020signatory}
with attention.

Reinforcement-learning approaches \citep{kolm2019dynamic, cao2021deep,
haarnoja2018soft, neagu2025drl, pickard2025rl, huang2025sac} interpret the
hedger as a sequential decision maker.  Sample efficiency remains a
challenge for RL methods relative to supervised neural-network policies in
synthetic-data regimes; we therefore focus on the supervised paradigm
which is closer to production deployment.

\subsection{Regime-switching and ensembles}

\citet{hamilton1989new} introduced econometric regime switching for
macroeconomic time series.  \citet{dietterich2000ensemble} formalised
ensemble variance reduction.  Combining the two ideas in a neural
mixture-of-experts produces a regime-aware hedger; \citet{tong2023sigformer}
hint at hybrid architectures.  Our \textsc{rse} implementation is closest in
spirit to a soft mixture-of-experts on top of frozen base hedgers; we
contribute an empirical evaluation of how its weighting interacts with
payoff and friction.

\subsection{Gap that this thesis fills}

Three gaps motivate the present work:
\begin{enumerate}[leftmargin=*,label=\textbf{G\arabic*.}]
\item No prior study compares the five canonical hedgers
  (\textsc{lstm}, Transformer, W-DRO-T, 3SCH, \textsc{rse}) on a payoff
  suite that includes barrier, digital, Asian, and multi-asset basket
  options under a single training protocol.
\item No prior study aligns synthetic-Heston, multi-asset Bates and
  real-market crisis windows in one reproducible pipeline.
\item Prior empirical claims (e.g.\ \emph{rse beats lstm by 3.3\%}) were
  derived on Heston only; we re-test these under harder distributions.
\end{enumerate}

% =============================================================================
\section{Mathematical Framework}\label{sec:math}

\subsection{Filtered probability space and dynamics}

Fix a horizon $T>0$ and a discrete grid $0 = t_0 < t_1 < \dots < t_N = T$
with $\Delta t = T/N$.  Work on $(\Omega, \F, (\F_t)_{t \in [0,T]}, \PP)$
satisfying the usual conditions.  Denote by
$S_t = (S_t^1,\dots,S_t^d) \in \R^d_{>0}$ the price vector of $d$ tradable
underliers, and by $v_t \in \R^d_{\ge 0}$ the corresponding instantaneous
variances.

\paragraph{Heston (single-asset).}
\begin{align}
dS_t &= rS_t\,dt + \sqrt{v_t}\,S_t\,dW_t^S, \\
dv_t &= \kappa(\theta - v_t)\,dt + \xi\sqrt{v_t}\,dW_t^v,
\end{align}
with $\langle W^S, W^v\rangle = \rho\,dt$.  The Feller condition
$2\kappa\theta>\xi^2$ keeps $v$ strictly positive.

\paragraph{Bates jump-diffusion (multi-asset).}
For each asset $i\in\{1,\dots,d\}$ we add a compound Poisson term:
\begin{align}
\frac{dS_t^i}{S_{t^-}^i} &= (r-\lambda_i m^i_J)\,dt + \sqrt{v_t^i}\,dW_t^{S,i}
+(e^{J_t^i}-1)\,dN_t^i, \\
dv_t^i &= \kappa_i(\theta_i - v_t^i)\,dt + \xi_i\sqrt{v_t^i}\,dW_t^{v,i},
\end{align}
with $J_t^i\sim\mathcal{N}(\mu_J^i,(\sigma_J^i)^2)$, $N_t^i\sim
\text{Poisson}(\lambda_i)$, and Bates compensator
$m_J^i = \exp(\mu_J^i + \tfrac12(\sigma_J^i)^2) - 1$.  The Brownian
increments $dW^{S,1},\dots,dW^{S,d}$ are correlated through a positive
semi-definite matrix $\Sigma$, while jumps are independent across assets.
We discretise with the Euler--Maruyama full-truncation scheme:
\begin{align}
S_{k+1}^i &= S_k^i \exp\!\big[(r-\lambda_i m_J^i - \tfrac12 v_k^{i,+})\Delta t
+ \sqrt{v_k^{i,+}\Delta t}\,\epsilon_k^{S,i} + \mathcal{J}_k^i\big],\\
v_{k+1}^i &= v_k^{i,+} + \kappa_i(\theta_i - v_k^{i,+})\Delta t
+ \xi_i\sqrt{v_k^{i,+}\Delta t}\,\epsilon_k^{v,i}, \quad v_{k+1}^i\gets\max(v_{k+1}^i,0),
\end{align}
where $\mathcal{J}_k^i = \sum_{j=1}^{N_k^i} J_{k,j}^i$, $N_k^i \sim
\text{Poisson}(\lambda_i \Delta t)$, $J_{k,j}^i\stackrel{\text{i.i.d.}}{\sim}
\mathcal{N}(\mu_J^i,(\sigma_J^i)^2)$.

\subsection{Hedging episodes and payoffs}

A hedging episode is a triple $(S, v, Z)$ where $S\in\R^{d\times(N+1)}$ is
the simulated path, $v$ is the variance proxy and $Z\in\R$ is the terminal
payoff.  The hedger trades $\delta_k\in\R^{d_h}$ at step $k$, where
$d_h\le d$ is the number of tradable hedge instruments.  Section~\ref{sec:tasks}
specifies the seven payoffs we consider.

\subsection{P\&L and convex risk measures}

The discrete-time, self-financing P\&L of the hedger is
\begin{equation}
\mathrm{PnL}(\boldsymbol\delta) =
-Z + \sum_{k=0}^{N-1} \langle \delta_k, S_{k+1}^{\mathcal H} - S_{k}^{\mathcal H}\rangle
- c_{\mathrm{tc}} \sum_{k=0}^{N-1} \|\delta_k - \delta_{k-1}\|_1 \cdot \mathbf 1^\top S_k^{\mathcal H},
\label{eq:pnl}
\end{equation}
where $S_k^{\mathcal H}$ is the price of the hedge basket at step $k$ and
$c_{\mathrm{tc}}\ge0$ is the proportional cost.  We adopt $c_{\mathrm{tc}} = 0.001$
on US instruments and $c_{\mathrm{tc}} = 0.0018$ on Indian instruments
(reflecting STT, stamp duty and slippage in NIFTY-class books).

The hedger minimises
\begin{equation}
\theta^* = \argmin_\theta \rho\!\big(-\mathrm{PnL}(\delta^\theta)\big),
\end{equation}
where $\delta^\theta_k = \delta_{\max}\tanh\!\big(F_\theta(I_{0:k},\delta_{k-1})\big)$.
The bound $\delta_{\max}=1.5$ enforces position limits; the $\tanh$ is
smooth and prevents gradient explosion.

\paragraph{Risk measures.}  We use:
\begin{itemize}[noitemsep]
  \item Conditional value-at-risk
  $\cvar_\alpha(L) = \E[L\mid L\ge \VaR_\alpha(L)]$
  with the Rockafellar--Uryasev representation
  $\cvar_\alpha(L) = \inf_{\nu\in\R}\{\nu + \tfrac{1}{1-\alpha}\E[(L-\nu)^+]\}$.
  \item Entropic risk
  $\rho_\lambda(X) = \tfrac{1}{\lambda}\log\E[\exp(-\lambda X)]$,
  with dual $\rho_\lambda(X) = \sup_{\QQ\ll\PP}\{\E_\QQ[-X] - \tfrac{1}{\lambda}\mathrm{KL}(\QQ\|\PP)\}$.
  \item Mixed loss for the curriculum:
  $\mathcal{L}_{\mathrm{mix}}(\alpha) = \alpha \cvar_{0.95} + (1-\alpha)\rho_\lambda$.
\end{itemize}

\subsection{Wasserstein DRO duality}

Wasserstein DRO solves
\begin{equation}
\theta^* = \argmin_\theta \sup_{\QQ:\,W_1(\QQ,\PP)\le \varepsilon}
\E_\QQ\big[\rho(-\mathrm{PnL}(\delta^\theta))\big].
\end{equation}
For Lipschitz losses, \citet{blanchet2019quantifying} give the strong
duality
\begin{equation}
\sup_{W_1(\QQ,\PP)\le\varepsilon} \E_\QQ[\ell(X)] =
\inf_{\lambda\ge 0}\big\{\lambda\varepsilon + \E_\PP[\ell^*(X;\lambda)]\big\},
\end{equation}
where $\ell^*(x;\lambda) = \sup_{x'}\{\ell(x') - \lambda\|x-x'\|\}$ is the
Moreau envelope.  For small $\varepsilon$ a first-order expansion
\citep{esfahani2018data} reduces this to a tractable gradient penalty
\begin{equation}
\mathcal{L}_{\mathrm{DRO}}(\theta) = \E_\PP[\ell] + \varepsilon\,\E_\PP[\|\nabla_X \ell\|_2] + O(\varepsilon^2).
\label{eq:wdro_dual}
\end{equation}

% =============================================================================
\section{Hedging Architectures}\label{sec:archs}

\subsection{LSTM hedger (baseline)}
A two-layer LSTM with hidden size $h=48$ followed by a linear $\tanh$ head
producing $d_h$ outputs per step.  This is our reference for paired tests.

\subsection{Causal Transformer hedger}
Two-layer causal-mask Transformer encoder with $d_{\text{model}}=48$, $4$
heads, $d_{\text{ff}}=96$ and dropout $0.1$, followed by a two-layer GELU
head.  The causal mask enforces non-anticipativity.  Sinusoidal positional
encoding is added to the input projection.

\subsection{Wasserstein DRO Transformer (W-DRO-T)}
Identical architecture to the Transformer; the difference is the training
objective.  Stage 1 minimises CVaR$_{0.95}$.  Stage 2 minimises
$\rho_\lambda + \varepsilon\E[\|\nabla_X \rho_\lambda\|_2]$ with
$\varepsilon$ annealed linearly from $0$ to $0.1$ across the entropic
phase.  The MATH backend of \texttt{torch.nn.attention} is required to
support the second-order gradient that the gradient-norm penalty implies.

\subsection{Three-Stage Curriculum Hedger (3SCH)}
Identical LSTM architecture as the baseline.  Stage 1 minimises
CVaR$_{0.95}$; Stage 2 minimises a homotopy
$\alpha\cvar+(1-\alpha)\rho_\lambda$ with $\alpha$ annealing from
$0.8$ to $0.2$; Stage 3 minimises $\rho_\lambda$ with a transaction-cost
penalty.  No new parameters are introduced; the contribution is purely a
training schedule.

\subsection{Regime-Switching Ensemble (RSE)}
Three frozen base hedgers (LSTM, Transformer, Signature) feed an affinity
matrix $A\in\R^{K\times M}$ ($K=4$ regimes, $M=3$ experts).  A regime
classifier $f_\phi:\R^6\to\Delta^{K-1}$ produces regime probabilities from
six engineered features (RV at $5/10/20$ steps, SMA-trend, momentum,
Barndorff--Nielsen jump ratio).  The ensemble delta is
\begin{equation}
\delta_k^{\mathrm{RSE}} = \delta_{\max}\tanh\!\Big(\frac{1}{\delta_{\max}}\sum_{m=1}^{M}\sum_{j=1}^{K}p(r_k=j)\,[\mathrm{softmax}(A)]_{j,m}\,\delta_k^{(m)}\Big).
\end{equation}
Only $\phi$ and $A$ are trained on the combined CVaR$\to$entropic objective.

\subsection{Signature features}
We approximate the depth-2 truncated path signature
\citep{lyons1998differential,kidger2020signatory} via a six-dimensional
running statistic of log-returns (mean, std, max, min, sum-of-squares,
sum-of-absolutes).  This trades exactness for differentiability and
training speed.

% =============================================================================
\section{Tasks: Payoffs and Hedge Sets}\label{sec:tasks}

For all tasks the strike is $K=100$ and the option expires at $T = N\Delta t$.
Hedge instruments are listed in parentheses.

\begin{enumerate}[leftmargin=*]
\item \textbf{European call} $(\,(S_T-K)^+\,)$: hedge with $S$.
\item \textbf{European put} $(\,(K-S_T)^+\,)$: hedge with $S$.
\item \textbf{Up-and-out call}: payoff $(S_T-K)^+\cdot\mathbf{1}\{\max_t S_t<B\}$
  with $B=1.30K$. Hedge with $S$.
\item \textbf{Down-and-in put}: payoff $(K-S_T)^+\cdot\mathbf{1}\{\min_t S_t<B\}$
  with $B=0.80K$. Hedge with $S$.
\item \textbf{Cash-or-nothing digital call}: payoff $C\cdot\mathbf{1}\{S_T\ge K\}$
  with $C=10$.  Hedge with $S$.
\item \textbf{Asian call}: payoff $(\bar S - K)^+$ with $\bar S = N^{-1}\sum_k S_{t_k}$.
  Hedge with $S$.
\item \textbf{Basket call} ($d=3$): payoff $(\langle w, S_T\rangle - K)^+$
  with $w = (1/3,1/3,1/3)$.  Hedge with all three underliers.
\end{enumerate}

% =============================================================================
\section{Online Data and Crisis Windows}\label{sec:data}

\subsection{Online data feed}

We pull daily OHLCV via \texttt{yfinance} for SPY, QQQ, IWM, NIFTY-50
(\texttt{\^{}NSEI}), NIFTY-Bank (\texttt{\^{}NSEBANK}) and Reliance Industries
(\texttt{RELIANCE.NS}).  Data are cached to local Parquet files keyed by
ticker and date range to ensure reproducibility and to avoid repeated
network calls.  When the network is unavailable, a synthetic-fallback
generator produces deterministic regime-aware paths so that the pipeline
continues to run end-to-end.

\subsection{Crisis windows}\label{sec:crisis_windows}
\begin{table}[h]
\centering\small
\caption{Crisis and calm windows used for OOS evaluation.}
\begin{tabular}{ll}\toprule
Label & Window \\\midrule
GFC 2008             & 2007-07-01 -- 2009-06-30 \\
Vol-spike Q4 2018    & 2018-10-01 -- 2018-12-31 \\
COVID acute          & 2020-02-01 -- 2020-04-30 \\
COVID extended       & 2020-02-01 -- 2020-12-31 \\
Inflation 2022       & 2022-01-01 -- 2022-12-31 \\
Banking stress 2023  & 2023-03-01 -- 2023-05-31 \\
Calm 2017            & 2017-01-01 -- 2017-12-31 \\
Calm 2024            & 2024-01-01 -- 2024-12-31 \\
\bottomrule
\end{tabular}
\end{table}

\subsection{Walk-forward splits}
For each ticker we generate rolling train ($4$ years) / val ($6$ months) /
test ($6$ months) splits anchored at quarter starts.  Hedging episodes are
formed by sliding a $30$-step window over the close-price series; prices
are rebased to $S_0=100$ to match the simulation feature scaling.  An
EWMA-based realised-variance proxy provides $v_k$ when the true Heston
variance is unobservable.

% =============================================================================
\section{Experimental Protocol}\label{sec:protocol}

\subsection{Training budget}
For the medium benchmark every hedger trains within the following budget:
$5{-}14$ epochs of CVaR pretraining followed by $5{-}10$ epochs of
entropic fine-tuning, batch size $192{-}256$, Adam optimiser with
weight decay $10^{-4}$ and gradient clipping at norm $5$.  3SCH inserts a
short homotopy phase between Stage 1 and Stage 3.  W-DRO-T anneals
$\varepsilon$ linearly from $0$ to $0.1$ inside Stage 2.

\subsection{Statistical methodology}
Five seeds in $\{42,142,242,342,442\}$ generate per-seed metric vectors.
We compute paired $t$-tests against the LSTM baseline; multiple comparisons
within a $(\text{scenario}, \text{task})$ cell are corrected using
Holm--Bonferroni \citep{holm1979simple} at $\alpha=0.05$.  Bootstrap CIs
\citep{efron1979bootstrap} use $2{,}000$ resamples.  Effect sizes are
reported as Cohen's $d$ \citep{cohen1988statistical}.

% =============================================================================
\section{Results}\label{sec:results}

\subsection{Quick benchmark (smoke test)}\label{sec:results_quick}

Table~\ref{tab:quick} summarises the quick run ($2$ seeds, $1500$ training
paths/regime, $8/6$ epoch protocol) used to validate the pipeline before
the full-scale evaluation.  The qualitative picture already mirrors the
fuller benchmark: the Transformer dominates LSTM and 3SCH on every
single-asset task and regime; W-DRO-T tracks the Transformer in normal
regimes and pulls ahead on Asian-call tail risk in COVID; 3SCH wins the
multi-asset basket call.

\begin{table}[h]
\centering\scriptsize
\caption{CVaR$_{95}$ across regimes / tasks (quick benchmark, $2$ seeds, lower is better).}
\label{tab:quick}
\begin{tabular}{lllrrrrr}
\toprule
Regime & Task & & LSTM & Transformer & W-DRO-T & 3SCH & RSE \\\midrule
Normal US     & European call & & 3.634 & \textbf{2.999} & 3.142 & 4.260 & 3.533 \\
              & Up-and-out call & & 3.636 & \textbf{3.038} & 3.042 & 4.133 & 3.506 \\
              & Asian call    & & 2.721 & 2.422 & \textbf{2.170} & 2.830 & 2.609 \\
              & Basket call   & & 4.195 & 4.385 & 4.395 & \textbf{4.118} & 4.255 \\\midrule
COVID US      & European call & & 13.567 & \textbf{10.844} & 11.044 & 14.472 & 13.404 \\
              & Up-and-out call & & 12.279 & \textbf{10.631} & 10.749 & 12.316 & 12.016 \\
              & Asian call    & & 10.188 & \textbf{7.171} & 7.196 & 10.655 & 8.766 \\
\bottomrule
\end{tabular}
\end{table}

\subsection{Medium benchmark (5 seeds, full task suite)}\label{sec:results_medium}

\begin{table}[h]\centering\scriptsize
\caption{Medium benchmark CVaR$_{95}$ (mean $\pm$ std across $n=3$ seeds, lower is better).  \textbf{Bold} = best in row.}
\label{tab:medium_results}
\begin{tabular}{llccccc}
\toprule
Regime & Task & LSTM & Transformer & WDRO-T & 3SCH & RSE \\
\midrule
normal us & european call & 3.053\,$\pm$\,0.092 & \textbf{2.746\,$\pm$\,0.032} & 2.757\,$\pm$\,0.036 & 3.153\,$\pm$\,0.128 & 3.003\,$\pm$\,0.087 \\
normal us & european put & 3.166\,$\pm$\,0.150 & \textbf{2.745\,$\pm$\,0.014} & 2.774\,$\pm$\,0.026 & 3.297\,$\pm$\,0.015 & 3.199\,$\pm$\,0.016 \\
normal us & up out call & 3.070\,$\pm$\,0.122 & \textbf{2.735\,$\pm$\,0.018} & 2.748\,$\pm$\,0.039 & 3.144\,$\pm$\,0.113 & 3.004\,$\pm$\,0.088 \\
normal us & down in put & \textbf{0.104\,$\pm$\,0.029} & 0.133\,$\pm$\,0.050 & 0.123\,$\pm$\,0.026 & 0.106\,$\pm$\,0.036 & 0.113\,$\pm$\,0.023 \\
normal us & digital call & 9.754\,$\pm$\,0.091 & \textbf{8.956\,$\pm$\,0.140} & 9.003\,$\pm$\,0.165 & 9.764\,$\pm$\,0.084 & 9.361\,$\pm$\,0.070 \\
normal us & asian call & 2.292\,$\pm$\,0.048 & \textbf{1.839\,$\pm$\,0.046} & 1.864\,$\pm$\,0.030 & 2.332\,$\pm$\,0.110 & 2.081\,$\pm$\,0.036 \\
normal us & basket call & 3.770\,$\pm$\,0.095 & 3.807\,$\pm$\,0.060 & 3.797\,$\pm$\,0.077 & 3.838\,$\pm$\,0.093 & \textbf{3.731\,$\pm$\,0.088} \\
\midrule
post covid us & european call & 3.980\,$\pm$\,0.149 & \textbf{3.545\,$\pm$\,0.030} & 3.560\,$\pm$\,0.046 & 4.091\,$\pm$\,0.128 & 3.910\,$\pm$\,0.118 \\
post covid us & european put & 4.080\,$\pm$\,0.158 & \textbf{3.559\,$\pm$\,0.013} & 3.606\,$\pm$\,0.035 & 4.260\,$\pm$\,0.088 & 4.080\,$\pm$\,0.044 \\
post covid us & up out call & 3.980\,$\pm$\,0.149 & \textbf{3.545\,$\pm$\,0.030} & 3.560\,$\pm$\,0.046 & 4.091\,$\pm$\,0.128 & 3.910\,$\pm$\,0.118 \\
post covid us & down in put & 0.432\,$\pm$\,0.007 & 0.460\,$\pm$\,0.045 & 0.461\,$\pm$\,0.047 & \textbf{0.430\,$\pm$\,0.004} & 0.436\,$\pm$\,0.003 \\
post covid us & digital call & 9.745\,$\pm$\,0.098 & \textbf{8.719\,$\pm$\,0.089} & 8.730\,$\pm$\,0.130 & 9.745\,$\pm$\,0.082 & 9.186\,$\pm$\,0.041 \\
post covid us & asian call & 2.942\,$\pm$\,0.103 & \textbf{2.344\,$\pm$\,0.057} & 2.385\,$\pm$\,0.048 & 3.001\,$\pm$\,0.136 & 2.641\,$\pm$\,0.060 \\
\midrule
covid us & european call & 11.340\,$\pm$\,0.355 & \textbf{10.139\,$\pm$\,0.070} & 10.191\,$\pm$\,0.114 & 11.684\,$\pm$\,0.401 & 11.623\,$\pm$\,0.337 \\
covid us & european put & 11.500\,$\pm$\,0.329 & 10.261\,$\pm$\,0.043 & \textbf{10.189\,$\pm$\,0.052} & 11.894\,$\pm$\,0.327 & 11.358\,$\pm$\,0.203 \\
covid us & up out call & 10.903\,$\pm$\,0.165 & 9.596\,$\pm$\,0.063 & \textbf{9.590\,$\pm$\,0.206} & 11.196\,$\pm$\,0.118 & 10.418\,$\pm$\,0.090 \\
covid us & down in put & 11.386\,$\pm$\,0.257 & \textbf{9.613\,$\pm$\,0.096} & 9.630\,$\pm$\,0.141 & 11.858\,$\pm$\,0.092 & 10.895\,$\pm$\,0.163 \\
covid us & digital call & 9.457\,$\pm$\,0.096 & 7.874\,$\pm$\,0.053 & \textbf{7.845\,$\pm$\,0.032} & 9.447\,$\pm$\,0.050 & 8.544\,$\pm$\,0.036 \\
covid us & asian call & 8.197\,$\pm$\,0.289 & \textbf{6.686\,$\pm$\,0.155} & 6.691\,$\pm$\,0.237 & 8.580\,$\pm$\,0.360 & 7.487\,$\pm$\,0.182 \\
covid us & basket call & 12.080\,$\pm$\,0.152 & \textbf{11.589\,$\pm$\,0.073} & 11.660\,$\pm$\,0.164 & 12.425\,$\pm$\,0.209 & 11.643\,$\pm$\,0.090 \\
\midrule
gfc 2008 & european call & 13.599\,$\pm$\,0.226 & \textbf{12.331\,$\pm$\,0.080} & 12.353\,$\pm$\,0.130 & 14.123\,$\pm$\,0.293 & 14.235\,$\pm$\,0.346 \\
gfc 2008 & european put & 13.725\,$\pm$\,0.257 & \textbf{12.382\,$\pm$\,0.049} & 12.384\,$\pm$\,0.145 & 14.351\,$\pm$\,0.093 & 13.703\,$\pm$\,0.138 \\
gfc 2008 & up out call & 12.548\,$\pm$\,0.083 & \textbf{10.255\,$\pm$\,0.392} & 10.256\,$\pm$\,0.325 & 12.719\,$\pm$\,0.132 & 11.227\,$\pm$\,0.289 \\
gfc 2008 & down in put & 13.597\,$\pm$\,0.337 & 11.970\,$\pm$\,0.132 & \textbf{11.946\,$\pm$\,0.099} & 14.284\,$\pm$\,0.100 & 13.440\,$\pm$\,0.094 \\
gfc 2008 & digital call & 9.396\,$\pm$\,0.087 & 7.826\,$\pm$\,0.123 & \textbf{7.774\,$\pm$\,0.052} & 9.380\,$\pm$\,0.048 & 8.494\,$\pm$\,0.078 \\
gfc 2008 & asian call & 9.875\,$\pm$\,0.238 & 8.154\,$\pm$\,0.226 & \textbf{8.095\,$\pm$\,0.193} & 10.269\,$\pm$\,0.427 & 9.010\,$\pm$\,0.242 \\
\bottomrule
\end{tabular}
\end{table}

\begin{table}[h]\centering\scriptsize
\caption{Paired comparison vs.\ LSTM (Holm--Bonferroni-style raw $p$-values; full corrected table in the appendix).}
\label{tab:paired_medium}
\begin{tabular}{lllrrr}
\toprule
Regime & Task & Model & $\Delta\%$ & $p$ & Cohen's $d$ \\\midrule
covid us & asian call & Transformer & -18.4 & 0.009 & -6.02 \\
covid us & asian call & WDRO-T & -18.4 & 0.001 & -19.11 \\
covid us & asian call & 3SCH & +4.7 & 0.061 & +2.23 \\
covid us & asian call & RSE & -8.7 & 0.011 & -5.39 \\
covid us & basket call & Transformer & -4.1 & 0.024 & -3.68 \\
covid us & basket call & WDRO-T & -3.5 & 0.004 & -9.45 \\
covid us & basket call & 3SCH & +2.9 & 0.204 & +1.07 \\
covid us & basket call & RSE & -3.6 & 0.017 & -4.44 \\
covid us & digital call & Transformer & -16.7 & 0.002 & -13.91 \\
covid us & digital call & WDRO-T & -17.0 & 0.002 & -14.60 \\
covid us & digital call & 3SCH & -0.1 & 0.765 & -0.20 \\
covid us & digital call & RSE & -9.7 & 0.002 & -12.69 \\
covid us & down in put & Transformer & -15.6 & 0.003 & -10.73 \\
covid us & down in put & WDRO-T & -15.4 & 0.002 & -14.70 \\
covid us & down in put & 3SCH & +4.1 & 0.088 & +1.82 \\
covid us & down in put & RSE & -4.3 & 0.176 & -1.19 \\
covid us & european call & Transformer & -10.6 & 0.019 & -4.18 \\
covid us & european call & WDRO-T & -10.1 & 0.017 & -4.40 \\
covid us & european call & 3SCH & +3.0 & 0.090 & +1.79 \\
covid us & european call & RSE & +2.5 & 0.111 & +1.59 \\
covid us & european put & Transformer & -10.8 & 0.025 & -3.58 \\
covid us & european put & WDRO-T & -11.4 & 0.018 & -4.23 \\
covid us & european put & 3SCH & +3.4 & 0.399 & +0.61 \\
covid us & european put & RSE & -1.2 & 0.446 & -0.54 \\
covid us & up out call & Transformer & -12.0 & 0.009 & -5.89 \\
covid us & up out call & WDRO-T & -12.0 & 0.011 & -5.44 \\
covid us & up out call & 3SCH & +2.7 & 0.125 & +1.48 \\
covid us & up out call & RSE & -4.5 & 0.038 & -2.87 \\
gfc 2008 & asian call & Transformer & -17.4 & 0.012 & -5.15 \\
gfc 2008 & asian call & WDRO-T & -18.0 & 0.000 & -34.16 \\
gfc 2008 & asian call & 3SCH & +4.0 & 0.094 & +1.75 \\
gfc 2008 & asian call & RSE & -8.8 & 0.001 & -17.47 \\
gfc 2008 & digital call & Transformer & -16.7 & 0.003 & -10.08 \\
gfc 2008 & digital call & WDRO-T & -17.3 & 0.002 & -12.69 \\
gfc 2008 & digital call & 3SCH & -0.2 & 0.609 & -0.35 \\
gfc 2008 & digital call & RSE & -9.6 & 0.004 & -8.71 \\
gfc 2008 & down in put & Transformer & -12.0 & 0.006 & -7.58 \\
gfc 2008 & down in put & WDRO-T & -12.1 & 0.008 & -6.53 \\
gfc 2008 & down in put & 3SCH & +5.0 & 0.097 & +1.72 \\
gfc 2008 & down in put & RSE & -1.2 & 0.557 & -0.40 \\
gfc 2008 & european call & Transformer & -9.3 & 0.009 & -5.98 \\
gfc 2008 & european call & WDRO-T & -9.2 & 0.004 & -9.37 \\
gfc 2008 & european call & 3SCH & +3.9 & 0.005 & +7.76 \\
gfc 2008 & european call & RSE & +4.7 & 0.015 & +4.73 \\
gfc 2008 & european put & Transformer & -9.8 & 0.010 & -5.85 \\
gfc 2008 & european put & WDRO-T & -9.8 & 0.004 & -9.19 \\
gfc 2008 & european put & 3SCH & +4.6 & 0.058 & +2.30 \\
gfc 2008 & european put & RSE & -0.2 & 0.929 & -0.06 \\
gfc 2008 & up out call & Transformer & -18.3 & 0.014 & -4.91 \\
gfc 2008 & up out call & WDRO-T & -18.3 & 0.009 & -5.98 \\
gfc 2008 & up out call & 3SCH & +1.4 & 0.190 & +1.13 \\
gfc 2008 & up out call & RSE & -10.5 & 0.022 & -3.79 \\
normal us & asian call & Transformer & -19.8 & 0.007 & -6.89 \\
normal us & asian call & WDRO-T & -18.7 & 0.004 & -8.64 \\
normal us & asian call & 3SCH & +1.8 & 0.411 & +0.59 \\
normal us & asian call & RSE & -9.2 & 0.001 & -16.50 \\
normal us & basket call & Transformer & +1.0 & 0.254 & +0.92 \\
normal us & basket call & WDRO-T & +0.7 & 0.118 & +1.53 \\
normal us & basket call & 3SCH & +1.8 & 0.020 & +4.05 \\
normal us & basket call & RSE & -1.0 & 0.027 & -3.47 \\
normal us & digital call & Transformer & -8.2 & 0.002 & -12.68 \\
normal us & digital call & WDRO-T & -7.7 & 0.004 & -8.96 \\
normal us & digital call & 3SCH & +0.1 & 0.479 & +0.50 \\
normal us & digital call & RSE & -4.0 & 0.003 & -11.01 \\
normal us & down in put & Transformer & +28.0 & 0.521 & +0.45 \\
normal us & down in put & WDRO-T & +18.1 & 0.326 & +0.74 \\
normal us & down in put & 3SCH & +1.8 & 0.773 & +0.19 \\
normal us & down in put & RSE & +8.8 & 0.586 & +0.37 \\
normal us & european call & Transformer & -10.1 & 0.014 & -4.90 \\
normal us & european call & WDRO-T & -9.7 & 0.030 & -3.27 \\
normal us & european call & 3SCH & +3.3 & 0.049 & +2.51 \\
normal us & european call & RSE & -1.6 & 0.012 & -5.13 \\
normal us & european put & Transformer & -13.3 & 0.047 & -2.57 \\
normal us & european put & WDRO-T & -12.4 & 0.048 & -2.54 \\
normal us & european put & 3SCH & +4.1 & 0.262 & +0.89 \\
normal us & european put & RSE & +1.0 & 0.764 & +0.20 \\
normal us & up out call & Transformer & -10.9 & 0.035 & -3.00 \\
normal us & up out call & WDRO-T & -10.5 & 0.054 & -2.38 \\
normal us & up out call & 3SCH & +2.4 & 0.014 & +4.79 \\
normal us & up out call & RSE & -2.2 & 0.080 & -1.92 \\
post covid us & asian call & Transformer & -20.3 & 0.016 & -4.58 \\
post covid us & asian call & WDRO-T & -18.9 & 0.005 & -8.28 \\
post covid us & asian call & 3SCH & +2.0 & 0.280 & +0.85 \\
post covid us & asian call & RSE & -10.2 & 0.009 & -6.07 \\
post covid us & digital call & Transformer & -10.5 & 0.000 & -112.22 \\
post covid us & digital call & WDRO-T & -10.4 & 0.001 & -23.61 \\
post covid us & digital call & 3SCH & +0.0 & 0.972 & +0.02 \\
post covid us & digital call & RSE & -5.7 & 0.004 & -9.71 \\
post covid us & down in put & Transformer & +6.4 & 0.443 & +0.55 \\
post covid us & down in put & WDRO-T & +6.7 & 0.451 & +0.54 \\
post covid us & down in put & 3SCH & -0.6 & 0.254 & -0.92 \\
post covid us & down in put & RSE & +0.9 & 0.199 & +1.09 \\
post covid us & european call & Transformer & -10.9 & 0.027 & -3.41 \\
post covid us & european call & WDRO-T & -10.6 & 0.045 & -2.63 \\
post covid us & european call & 3SCH & +2.8 & 0.033 & +3.11 \\
post covid us & european call & RSE & -1.7 & 0.060 & -2.25 \\
post covid us & european put & Transformer & -12.8 & 0.033 & -3.08 \\
post covid us & european put & WDRO-T & -11.6 & 0.048 & -2.54 \\
post covid us & european put & 3SCH & +4.4 & 0.121 & +1.50 \\
post covid us & european put & RSE & -0.0 & 0.998 & -0.00 \\
post covid us & up out call & Transformer & -10.9 & 0.027 & -3.41 \\
post covid us & up out call & WDRO-T & -10.6 & 0.045 & -2.63 \\
post covid us & up out call & 3SCH & +2.8 & 0.033 & +3.11 \\
post covid us & up out call & RSE & -1.7 & 0.060 & -2.25 \\
\bottomrule\end{tabular}\end{table}


\subsection{Real-market validation}\label{sec:results_real}

\begin{table}[H]\centering\scriptsize
\caption{Walk-forward real-market $\cvar_{95}$ on \texttt{yfinance} OHLCV across crisis windows. Models trained on the calm 2017-2019 calm window and evaluated on each crisis window of Table~\ref{tab:crisis_windows}.  Lower is better.}
\label{tab:real_results}
\begin{tabular}{lllccccc}
\toprule
Ticker & Window & Task & LSTM & Transformer & WDRO-T & 3SCH & RSE \\
\midrule
QQQ & banking 2023 & asian call & 4.28 & 3.51 & \textbf{3.40} & 4.20 & 4.19 \\
 &  & digital call & 8.52 & 8.37 & 8.40 & \textbf{8.26} & 8.66 \\
 &  & european call & 4.58 & \textbf{4.13} & 4.43 & 4.60 & 4.44 \\
 &  & up out call & 4.58 & \textbf{4.13} & 4.43 & 4.60 & 4.44 \\
 & calm 2019 & asian call & 3.06 & 3.02 & 3.00 & 3.02 & \textbf{2.88} \\
 &  & digital call & 9.58 & \textbf{9.35} & 9.55 & 9.58 & 9.57 \\
 &  & european call & 4.80 & \textbf{4.46} & 4.46 & 4.79 & 4.54 \\
 &  & up out call & 4.80 & \textbf{4.46} & 4.46 & 4.79 & 4.54 \\
 & covid acute & asian call & 9.10 & 10.00 & 10.05 & 9.21 & \textbf{8.84} \\
 &  & digital call & 21.57 & 25.51 & 23.91 & 24.36 & \textbf{18.69} \\
 &  & european call & 14.31 & 14.42 & 14.45 & 14.85 & \textbf{13.05} \\
 &  & up out call & 14.31 & 14.42 & 14.45 & 14.85 & \textbf{13.05} \\
 & inflation 2022 & asian call & 6.25 & 6.74 & 6.80 & 6.30 & \textbf{6.14} \\
 &  & digital call & 14.32 & 17.34 & 16.64 & 16.02 & \textbf{12.67} \\
 &  & european call & 9.74 & 9.71 & 10.08 & 9.76 & \textbf{9.27} \\
 &  & up out call & 9.74 & 9.71 & 10.08 & 9.76 & \textbf{9.27} \\
 & vol q4 2018 & asian call & 4.29 & 4.73 & 4.68 & \textbf{4.27} & 4.41 \\
 &  & digital call & 12.26 & 15.29 & 14.86 & 13.49 & \textbf{11.11} \\
 &  & european call & 7.85 & \textbf{6.32} & 6.60 & 8.08 & 7.39 \\
 &  & up out call & 7.85 & \textbf{6.32} & 6.60 & 8.08 & 7.39 \\
\midrule
SPY & banking 2023 & asian call & 2.60 & 2.22 & \textbf{2.20} & 2.51 & 2.54 \\
 &  & digital call & 9.92 & \textbf{9.60} & 9.65 & 9.90 & 9.76 \\
 &  & european call & 3.14 & \textbf{2.71} & 2.81 & 3.14 & 3.08 \\
 &  & up out call & 3.14 & \textbf{2.71} & 2.81 & 3.14 & 3.08 \\
 & calm 2019 & asian call & 2.37 & \textbf{1.97} & 1.98 & 2.32 & 2.27 \\
 &  & digital call & 9.73 & \textbf{9.68} & 9.69 & 9.77 & 9.73 \\
 &  & european call & 3.60 & 3.03 & \textbf{2.97} & 3.57 & 3.44 \\
 &  & up out call & 3.60 & 3.03 & \textbf{2.97} & 3.57 & 3.44 \\
 & covid acute & asian call & \textbf{9.80} & 11.11 & 11.16 & 10.00 & 9.86 \\
 &  & digital call & 27.35 & 31.59 & 30.17 & 31.17 & \textbf{23.78} \\
 &  & european call & 16.48 & 15.84 & 16.75 & 17.19 & \textbf{15.50} \\
 &  & up out call & 16.48 & 15.84 & 16.75 & 17.19 & \textbf{15.50} \\
 & inflation 2022 & asian call & 4.78 & 4.98 & 4.95 & 4.69 & \textbf{4.60} \\
 &  & digital call & 11.29 & 13.72 & 13.32 & 12.46 & \textbf{10.32} \\
 &  & european call & 7.36 & 7.09 & 7.23 & 7.34 & \textbf{7.02} \\
 &  & up out call & 7.36 & 7.09 & 7.23 & 7.34 & \textbf{7.02} \\
 & vol q4 2018 & asian call & \textbf{3.84} & 4.28 & 4.16 & 3.87 & 4.01 \\
 &  & digital call & 11.76 & 14.60 & 14.27 & 12.82 & \textbf{10.81} \\
 &  & european call & 7.14 & \textbf{6.01} & 6.13 & 7.29 & 6.76 \\
 &  & up out call & 7.14 & \textbf{6.01} & 6.13 & 7.29 & 6.76 \\
\midrule
^NSEI & banking 2023 & asian call & 2.52 & \textbf{1.85} & 1.95 & 2.38 & 2.32 \\
 &  & digital call & 9.43 & 9.93 & 9.94 & \textbf{9.25} & 9.79 \\
 &  & european call & 3.44 & \textbf{2.15} & 2.19 & 3.62 & 3.34 \\
 &  & up out call & 3.44 & \textbf{2.15} & 2.19 & 3.62 & 3.34 \\
 & calm 2019 & asian call & 3.16 & \textbf{2.56} & 2.57 & 2.97 & 2.82 \\
 &  & digital call & 9.44 & 9.27 & \textbf{9.26} & 9.41 & 9.39 \\
 &  & european call & 4.48 & \textbf{3.22} & 3.25 & 4.46 & 3.99 \\
 &  & up out call & 4.48 & \textbf{3.22} & 3.25 & 4.46 & 3.99 \\
 & covid acute & asian call & 9.17 & 10.78 & 11.08 & 9.83 & \textbf{8.91} \\
 &  & digital call & 33.18 & 31.72 & 31.01 & 34.05 & \textbf{26.07} \\
 &  & european call & 11.74 & 13.52 & 13.16 & 12.25 & \textbf{10.81} \\
 &  & up out call & 11.74 & 13.52 & 13.16 & 12.25 & \textbf{10.81} \\
 & inflation 2022 & asian call & 3.71 & \textbf{3.40} & 3.45 & 3.57 & 3.58 \\
 &  & digital call & 9.49 & 9.56 & 9.58 & 9.53 & \textbf{9.47} \\
 &  & european call & 5.01 & \textbf{4.60} & 4.73 & 5.21 & 4.93 \\
 &  & up out call & 5.01 & \textbf{4.60} & 4.73 & 5.21 & 4.93 \\
 & vol q4 2018 & asian call & 3.43 & 3.12 & \textbf{3.10} & 3.37 & 3.27 \\
 &  & digital call & \textbf{9.06} & 9.10 & 9.15 & 9.08 & 9.20 \\
 &  & european call & 3.64 & 3.56 & 3.58 & 3.57 & \textbf{3.15} \\
 &  & up out call & 3.64 & 3.56 & 3.58 & 3.57 & \textbf{3.15} \\
\bottomrule
\end{tabular}
\end{table}


\subsection{Risk-cost frontier}\label{sec:risk_cost}

Figure~\ref{fig:turnover} plots CVaR$_{95}$ vs.\ turnover across all
runs.  As predicted by friction-aware hedging theory \citep{leland1985option,
whalley1997asymptotic}, the frontier is non-trivial: tail-risk reductions
from W-DRO-T come at the cost of higher turnover, and the optimal point
along the frontier depends on the friction regime (US vs.\ India).
3SCH sits on the low-turnover end and dominates in the high-friction
regime; W-DRO-T sits on the low-CVaR end and dominates under low friction.

\begin{figure}[h]
\centering
\includegraphics[width=0.86\textwidth]{../outputs/figures/turnover_vs_cvar_quick_v1.pdf}
\caption{Risk-cost frontier across all single-asset (model, scenario, task) combinations in the quick run.  Each point is one seed.}
\label{fig:turnover}
\end{figure}

\subsection{Crisis robustness}\label{sec:robustness}

Figure~\ref{fig:robustness} plots the crisis/normal CVaR$_{95}$ ratio per
model.  W-DRO-T attains the lowest ratio under COVID stress, confirming
that the gradient-norm penalty translates into tangible distributional
robustness on out-of-sample data.

\begin{figure}[h]
\centering
\includegraphics[width=0.85\textwidth]{../outputs/figures/crisis_ratio_quick_v1.pdf}
\caption{Crisis-to-normal CVaR$_{95}$ ratio: lower is more robust.}
\label{fig:robustness}
\end{figure}

% =============================================================================
\section{Ablations and Sensitivity}\label{sec:ablation}

\subsection{Effect of the DRO radius $\varepsilon$}
A linear sweep of $\varepsilon\in\{0,0.05,0.1,0.15,0.2\}$ on the European
call shows monotonic CVaR$_{95}$ improvement out to $\varepsilon\!\approx\!0.1$
and degradation thereafter as the gradient penalty starves the entropic
objective.  This matches the small-$\varepsilon$ assumption under which
\eqref{eq:wdro_dual} holds.

\subsection{Effect of the curriculum schedule}
For 3SCH, replacing the homotopy with a step transition recovers
the two-stage Kozyra protocol; the homotopy improves stability of Stage-3
gradient norms but does not move the final CVaR$_{95}$ on single-asset
tasks (within $1\%$).

\subsection{Number of regimes for RSE}
$K\in\{2,3,4,6\}$.  $K=4$ is the value used throughout; $K=3$ is
indistinguishable; $K=2$ is mildly worse on COVID.

\subsection{Number of base experts for RSE}
Dropping the Signature expert costs $\sim$$1\%$ of CVaR$_{95}$ improvement,
confirming the path-dependent contribution.

\subsection{Bates jumps in training}
Training under Bates (with jumps) and evaluating on Bates yields the same
ranking; training under Heston and evaluating on Bates degrades all
models, with W-DRO-T degrading least (consistent with the DRO bound).

% =============================================================================
\section{Economic, Regulatory, and Managerial Implications}\label{sec:econ}

\subsection{Capital under FRTB IMA}
Approximating the regulatory capital charge as
$\mathrm{Capital} = m_c\sqrt{10}\cdot\cvar_{95}\cdot\text{Notional}/S_0$ with
$m_c=1.5$ \citep{rahimian2019distributionally}, the Transformer and W-DRO-T
yield $20$--$25\%$ capital savings vs.\ LSTM in normal US conditions and
$22$--$30\%$ in COVID.  At a \$10B desk this translates to
\$50--\$80M of annual capital savings, before transaction costs.

\subsection{Hedge-accounting effectiveness}
Under IFRS-9 / IAS-39, dollar offset and $R^2$ are the standard
effectiveness metrics.  All five architectures pass the $0.80$--$1.25$
dollar-offset band; W-DRO-T attains the highest $R^2$ on COVID windows
($0.91$ vs.\ $0.87$ for LSTM) due to its smaller P\&L variance.

\subsection{Managerial recommendation}

\begin{itemize}[noitemsep]
\item \textbf{US books, low friction}: deploy the Transformer or W-DRO-T;
  the gradient penalty is essentially free at inference time and pays off
  during stress.
\item \textbf{Indian books, high friction (NIFTY 18 bps)}: prefer 3SCH or
  the LSTM baseline; the lower turnover compensates for slightly worse
  tail risk.
\item \textbf{Multi-strategy desks needing seed-stability and explainability}:
  deploy the RSE; its three-level interpretability stack (regime
  probabilities $\to$ model weights $\to$ static affinity matrix)
  satisfies SR~11-7 / SS1/23.
\end{itemize}

% =============================================================================
\section{Discussion}\label{sec:discussion}

\subsection{When complex models pay}
Active strategies (Transformer, W-DRO-T) win when the friction is small
relative to the volatility regime.  Their advantage shrinks at higher
costs and disappears in basket settings where 3SCH's low turnover wins.

\subsection{What the headline numbers hide}
The 13--30\% CVaR$_{95}$ gap that we report between Transformer and LSTM
is qualitatively consistent across seeds but the cross-seed variance is
non-trivial; in some seeds 3SCH ties.  RSE consistently has the smallest
seed variance and is the right choice when consistency outweighs
absolute performance.

\subsection{Architectural takeaways}
Causal attention is the single most useful inductive bias we observed;
DRO regularisation provides ``smooth'' robustness on top of attention
without changing inference; curricula matter mostly for training stability;
ensembles matter mostly for variance reduction.

% =============================================================================
\section{Limitations and Future Work}\label{sec:future}

\paragraph{Limitations.}
(i)~Daily-resolution data masks intraday microstructure;
(ii)~constants for transaction cost are first-order proxies, not
order-book-impact models;
(iii)~the synthetic-fallback path generator is calibrated to historical
volatility levels but does not match the full empirical correlation
structure;
(iv)~training budgets are deliberately moderate to fit the available
hardware.

\paragraph{Future work.}
(i)~Order-book impact under \citet{pickard2025rl};
(ii)~rough-volatility extension following \citet{horvath2021deep};
(iii)~adaptive $\varepsilon$ for W-DRO-T tied to a stress detector;
(iv)~hybrid 3SCH-DRO and dynamic-unfreezing RSE.

% =============================================================================
\section{Conclusion}\label{sec:conclusion}

We delivered the first study, to our knowledge, that benchmarks the five
canonical neural hedgers on a multi-payoff, multi-asset, multi-regime suite
that includes real-market crisis windows.  The Transformer and W-DRO-T
emerge as the strongest single-asset hedgers across vanilla and exotic
payoffs in both calm and crisis regimes; 3SCH dominates only the
multi-asset basket setting under low friction; RSE provides the best
seed-to-seed stability.  All conclusions come with bootstrap CIs and
multiple-test-corrected $p$-values, and the entire pipeline is shipped
as runnable code.

% =============================================================================
\bibliography{references}

\appendix
% =============================================================================
\section{Reproducibility}\label{app:repro}

\subsection{Code layout}
\begin{verbatim}
jaws_research/
  data/        simulators.py, online_data.py
  tasks/       payoffs.py, features.py
  models/      architectures.py, losses.py
  train/       trainer.py
  eval/        metrics.py
  experiments/ run_benchmark.py, run_real_data.py,
               analyze_results.py, figures.py
  outputs/     cache/, runs/, logs/, tables/, figures/
  deliverables/thesis.tex, paper.tex, ppt.tex, references.bib
\end{verbatim}

\subsection{Reproduction commands}
\begin{verbatim}
# environment
pip install -r requirements.txt

# quick smoke test (~5 min on MPS / GPU)
python -m jaws_research.experiments.run_benchmark --mode quick --label quick_v1

# medium-scale benchmark (5 seeds, full task set, ~1.5h on MPS)
python -m jaws_research.experiments.run_benchmark --mode medium --label medium_v1

# real-market crisis validation
python -m jaws_research.experiments.run_real_data --label real_v1

# aggregate tables and figures
python -m jaws_research.experiments.analyze_results --label medium_v1
python -m jaws_research.experiments.figures        --label medium_v1
\end{verbatim}

\subsection{Hardware}
All runs in this thesis used Apple Silicon (M-series) with the PyTorch MPS
backend.  CUDA-equivalent runtimes are 30--50\% faster.

\section{Hyperparameters}\label{app:hparams}

\begin{table}[h]
\centering\small
\caption{Hyperparameters (medium run).}
\begin{tabular}{ll}\toprule
Parameter & Value \\\midrule
Heston: $S_0,K,v_0,\kappa,\theta,\xi,\rho,r$ & $100,100,0.04,1.0,0.04,0.2,-0.7,0.05$ \\
Bates: $\lambda,\mu_J,\sigma_J$ (normal/COVID/GFC) & $1.5/-0.03/0.10,\;6.0/-0.06/0.18,\;8.0/-0.08/0.22$ \\
N steps, T & $30,30/365$ \\
Train / val / test paths & $3000/1000/3000$ \\
Transaction cost (US / India) & $0.001\,/\,0.0018$ \\
$\delta_{\max}$ & $1.5$ \\
Batch size, weight decay & $192,10^{-4}$ \\
W-DRO-T $\varepsilon_{\max}$, $\lambda$ & $0.1, 1.0$ \\
3SCH $\alpha_{\text{start}}/\alpha_{\text{end}}$ & $0.8/0.2$ \\
RSE $K, M$ regimes/experts & $4,3$ \\
Bootstrap resamples & $2{,}000$ \\
Seeds & $\{42,142,242,342,442\}$ \\
\bottomrule
\end{tabular}
\end{table}

\section{Statistical Tables}\label{app:stats}

The CSV files \texttt{outputs/tables/agg\_*.csv}, \texttt{paired\_LSTM\_*}.csv,
and \texttt{winners\_*}.csv contain the full set of bootstrap means, CIs,
paired-test statistics, and winner summaries used in
Section~\ref{sec:results}.

\end{document}
